\documentclass[12pt]{article}

\usepackage[margin=1in]{geometry}
\usepackage{setspace}
\usepackage{enumitem}
\usepackage{titlesec}

\titleformat{\section}{\large\bfseries}{\thesection}{1em}{}
\setstretch{1.2}

\title{UMBC Parking Management System\\Database Design Proposal}
\author{Mateo Jacome, Jason Chen}
\date{Mar 3 2026}

\begin{document}

\maketitle

\section*{Executive Summary}

\subsection*{Problem Statement}

Parking at a large commuter campus such as UMBC is inefficient, stressful, and difficult to enforce. Students frequently waste valuable time searching for open spaces. Visitors struggle to locate valid parking without clear guidance. Faculty and staff require access to restricted or premium areas. Meanwhile, enforcement is largely reactive and labor-intensive, relying on officers instead of automated validation mechanisms.

The core problem is the absence of a centralized, real-time database system that integrates parking permits, lot and spot management, vehicle data, reservations, sensor occupancy detection, and automated enforcement into a unified platform. Without such integration, parking operations lack efficiency, transparency, and scalability.

The proposed UMBC Parking Management System addresses this issue by implementing a PostgreSQL-backed database that manages parking as a real-time, data-driven ecosystem.

\subsection*{Parking Users}

The system supports multiple categories of users:

\begin{itemize}
    \item \textbf{Commuter Students} -- Purchase student permits and park in designated student lots.
    \item \textbf{Faculty and Staff} -- Hold faculty/staff permits and access restricted premium zones.
    \item \textbf{Visitors and Guest Lecturers} -- Reserve temporary parking spaces online and pay per visit.
    \item \textbf{Parking Administration} -- Manage permits, pricing structures, enforcement rules, and reporting.
    \item \textbf{Automated Enforcement System} -- Generates violations based on database validation rules.
    \item \textbf{Ground Sensors (IoT Devices)} -- Automatically update occupancy and trigger validation checks.
\end{itemize}

\subsection*{Scope of Implementation}

This project will implement a PostgreSQL database system that includes:

\begin{enumerate}
    \item \textbf{User and Vehicle Management}
    \begin{itemize}
        \item Store user profiles (students, faculty, visitors).
        \item Associate multiple vehicles with individual users.
        \item Store license plate information for validation.
    \end{itemize}

    \item \textbf{Permit Management}
    \begin{itemize}
        \item Define permit types (student, faculty, visitor).
        \item Assign permits to users.
        \item Enforce expiration dates.
        \item Restrict lot and zone access by permit type.
    \end{itemize}

    \item \textbf{Lot and Spot Management}
    \begin{itemize}
        \item Define parking lots, zones, rows, and individual spots.
        \item Track real-time occupancy per spot.
        \item Associate spots with allowed permit types.
    \end{itemize}

    \item \textbf{Reservation System}
    \begin{itemize}
        \item Allow visitors to reserve specific spots.
        \item Prevent double-booking through constraints and transactions.
        \item Lock spots upon confirmed payment.
    \end{itemize}

    \item \textbf{Sensor Integration}
    \begin{itemize}
        \item Update occupancy when a ground sensor detects a vehicle.
        \item Validate vehicle authorization against zone restrictions.
    \end{itemize}

    \item \textbf{Automated Ticketing}
    \begin{itemize}
        \item Generate violations for unauthorized parking.
        \item Store fine amounts and violation details.
        \item Record notification events.
    \end{itemize}

    \item \textbf{Payments}
    \begin{itemize}
        \item Record reservation payments.
        \item Record ticket payments.
        \item Track outstanding balances.
    \end{itemize}

    \item \textbf{Reporting and Analytics}
    \begin{itemize}
        \item Lot utilization rates.
        \item Revenue reports.
        \item Violation frequency reports.
        \item Active and expired permit tracking.
        \item Peak usage analysis.
    \end{itemize}

    \item \textbf{Concurrency Control}
    \begin{itemize}
        \item Prevent race conditions during reservation.
        \item Maintain consistency during simultaneous sensor updates.
    \end{itemize}
\end{enumerate}

This scope focuses specifically on the database design, implementation, constraints, indexing, realistic operations, and concurrency control. It does not include a full mobile or web application interface.

\section*{Concrete Use Cases}

The database must support the following operations:

\begin{itemize}
    \item Register a new user (student, faculty, or visitor).
    \item Add one or more vehicles to a user account.
    \item Purchase and assign a parking permit to a user.
    \item Define parking lots, zones, rows, and individual parking spots.
    \item Associate permit types with allowed zones.
    \item Query real-time available parking spots in a specific lot.
    \item Create a visitor reservation for a specific date and time.
    \item Prevent double-booking of a parking spot.
    \item Update occupancy when a sensor detects a vehicle.
    \item Validate a vehicle's permit against the spot's allowed zone.
    \item Automatically generate a ticket for unauthorized parking.
    \item Record ticket payments and mark violations as resolved.
    \item Generate reports on lot occupancy over time.
    \item Generate revenue reports from permits, reservations, and fines.
    \item Detect expired permits and flag vehicles as invalid.
    \item Handle simultaneous reservation attempts without data corruption.
\end{itemize}

\end{document}